%%%%%%%%%%%%%%%%%%%%%%%%%%%%%%%%%%%%%%%%%%%%%%%%%%%%%%%%%%%%%%%%%%%%%%%%%%%%%%
% COMP 9200 Assignment 3, Part 2 -- first draft of the proposal's
% methodology section (individual).
%
% Proposal: Beyond Complexity Reduction: Probability Sensitivity in
% LLM-Based Code Generation Systems (see assignment1/submission/part2).
%
% This file is a skeleton only: headings plus the questions each section
% has to answer. Replace every \todo{} with your own text.
% Build: latexmk -pdf COMP9200_A3_Part2_Methodology.tex
%%%%%%%%%%%%%%%%%%%%%%%%%%%%%%%%%%%%%%%%%%%%%%%%%%%%%%%%%%%%%%%%%%%%%%%%%%%%%%

\documentclass[conference]{IEEEtran}
\IEEEoverridecommandlockouts

\usepackage{cite}
\usepackage{amsmath,amssymb,amsfonts}
\usepackage{graphicx}
\usepackage{textcomp}
\usepackage{xcolor}
\usepackage{booktabs}
\usepackage[hidelinks]{hyperref}
\usepackage{url}

\newcommand{\todo}[1]{\textcolor{red}{[TODO: #1]}}

\title{Beyond Complexity Reduction: Probability Sensitivity in LLM-Based
Code Generation Systems\\[0.2em] \large Methodology (First Draft)}

\author{
\IEEEauthorblockN{Christopher Elliott}
\IEEEauthorblockA{\textit{COMP 9200: Applied Research Methods} \\
\textit{School of Computing and Academic Studies} \\
\textit{British Columbia Institute of Technology}, Burnaby, Canada \\
celliott64@my.bcit.ca}
}

\begin{document}
\maketitle

\section{Research Design}
\todo{Name the design and justify it. Which parts are experimental
(manipulated surface typicality, Subproblem 2) and which are
correlational or descriptive (Subproblems 1 and 3)? What are the
independent, dependent and control variables? State the hypotheses for
each subproblem. Relate the design to the controlled approach of McCoy et
al.\ \cite{mccoyEmbersAutoregressionUnderstanding2023}.}

\section{Data Collection}

\subsection{Tasks and Materials}
\todo{Which benchmark or task source (single-function / single-file only,
per your delimitations)? How many tasks, and how are they chosen? Why is
that number enough (power analysis or precedent)?}

\subsection{Constructing Matched Request Pairs}
\todo{How a ``typical'' and an ``unusual'' wording are written for each
task while keeping behaviour, test suite, language, library and algorithm
fixed. Who writes them, using what written guide, and how the second
person checks each pair.}

\subsection{Models and Scoring}
\todo{Which open-weight model(s) score input and output probability, and
which model(s) generate code. Exact versions, decoding settings
(temperature, samples per task), hardware. Total and per-token
log-probability, and why both.}

\subsection{Harness Condition (Subproblem 3)}
\todo{Which harness (e.g.\ Agentless), how its sub-prompts are captured
and scored, and what the single-prompt comparison is.}

\subsection{Procedure}
\todo{Step-by-step order of operations, so another researcher could repeat
it. What is logged at each step.}

\section{Data Analysis}
\todo{For each subproblem: the statistical test or model, the unit of
analysis, how repeated measures on the same task are handled (e.g.\ paired
tests or mixed-effects models), effect sizes, and correction for multiple
comparisons. How the two outcome measures (passes tests; matches
reference) are analysed. What result would support or refute each
hypothesis.}

\section{Validity and Reliability}

\subsection{Validity of Instruments}
\todo{Construct validity: does model log-probability measure
``typicality''? How Subproblem 1 checks this. Content validity of the
request pairs: how you show a pair differs only in surface form.}

\subsection{Reliability of Instruments}
\todo{Inter-rater agreement on pair equivalence (e.g.\ Cohen's $\kappa$,
with a threshold and what happens below it). Test--retest stability of
probability scores and of generated-code pass rates across seeds.}

\subsection{Threats to Internal and External Validity}
\todo{Confounds such as prompt length, benchmark contamination, and the
open-weight model standing in for closed ones. How each is controlled or
measured. What the results do and do not generalise to.}

\section{Ethical Considerations}
\todo{Licences of benchmarks and models, compute cost, whether any human
participants (the second rater) need consent, data and code release.}

\section{Limitations}
\todo{What the design cannot show, beyond the delimitations already
stated in the proposal.}

\appendices
\section{Use of AI Tools}
\todo{Your AI use statement for this part.}

\bibliographystyle{IEEEtran}
\bibliography{c9200references}

\end{document}
